% 원제: 14. 소자 수준 회로 모델(Simulink) 현황
\chapter{소자 수준 회로 모델(Simulink) 현황}

\section{테스트벤치 시리즈}

소자 수준 회로 모델은 하나의 파일을 계속 고치지 않고 검증 단계마다 새 테스트벤치를 만들어 남겼다. 모든 모델은 빌드 스크립트로 처음부터 생성되며, 단별 서브시스템(전력단·게이트·제어·계측), 역할이 드러나는 소자 이름, 결선·파라미터를 설명하는 주석 블록, 한글 주석이 달린 제어 코드를 갖춘다. 각 테스트벤치는 같은 조건의 검증 모델 결과를 기준값으로 두어 시험 전에 고정한 기준으로 판정한다(표 \ref{t:tb}). 10월 25일 보고 이후 TB14–TB20의 일곱 개를 더했다.

{\footnotesize
\begin{longtable}{L{\dimexpr 0.0965\linewidth-2\tabcolsep\relax}L{\dimexpr 0.3860\linewidth-2\tabcolsep\relax}L{\dimexpr 0.1754\linewidth-2\tabcolsep\relax}L{\dimexpr 0.3421\linewidth-2\tabcolsep\relax}}
\caption{테스트벤치 시리즈와 판정}\label{t:tb}\\
\toprule
\tabhead \textbf{TB} & \textbf{대상} & \textbf{기준} & \textbf{판정} \\
\midrule
\endfirsthead
\toprule
\tabhead \textbf{TB} & \textbf{대상} & \textbf{기준} & \textbf{판정} \\
\midrule
\endhead
\bottomrule
\endlastfoot
TB01 & Stage A 1위상 RSC, 고정 타이밍 & E1 & 통과(전 지표 0.2 \% 이내) \\
TB02 & Stage B H-브리지(개루프), 데드타임 전압오차 & 교정표 & 통과(최대 차 0.135 V) \\
TB03 & Stage B DAB(개루프), 위상 계단·직류 편류 & E12 & 통과(0.39 A 이내) \\
TB04 & Stage B 폐루프 & E12 & 일치 통과(리플은 TB04B에서 결론) \\
TB04B & TB04 + 후보 소자 출력용량, 게이트 격자 & E12 & 통과(0.625 ns, 리플 0.069·0.031 V) \\
TB05 & Stage C CLLC 3 kW(고정이득) & E10 & 통과(이득 0.9965–1.0001) \\
TB06 & 통합(주행): Stage A 모듈 1 + B + C & E12, E1, TB05 & 통과 3/3 \\
TB07 & 통합(충전): Stage A 8위상 + 충전기 + B + C & E16, E1 & 회로 한계 7/7, 일치는 조건부(충전전류 램프 ≥ 100 µs) \\
TB08 & 창 조정 모듈 R2-W460 + 바이패스·모드 시퀀서 & E12, E19 & 통과 5/5 \\
TB09 & R2 보호 상태기계, 제동 저항, 소프트 투입, 감쇠기 & E17 & 최종 실행 4/4(첫 실행 2사례 불통과 → 상태기계 보완) \\
TB10 & DAB 고장 모드 변조(EPS 게이트, 정합 유지, 감발) & E20 & 통과 3/3(회로 기준 2배), 1.5배 목표는 붕괴 구간 초과 \\
TB11 & TB10 + 데드타임 반영 첨두 표 & E21 & 불통과(DAB 첨두 88.5·68.9 A) — 지령표 칸 전환 과도 \\
TB12 & TB11 + EPS 지령 전이의 직류 편류 방지·변화율 상한 & E21 & 1/2 통과(1.5배 기준), DAB 첨두 70.0·68.5 A — 전이 끝 1.3 A 초과 \\
TB13 & TB12 + 지령 전이 여유(설계 상한 1.387배) & E21 & 통과 2/2(1.5배 기준), DAB 첨두 66.4·64.4 A \\
TB14 & 조정형 Stage C 5 kW(동기 벅 + CLLC 직류변압기 4.8:1) & E22 & 1차 0/5(영전압 스위칭) → 재설계 뒤 2차 6/7(270 V 역방향 경부하 불통과) \\
TB15 & TB14 + 경부하 두 브리지 구동 & E22 & 통과 5/5 \\
TB16 & R2 기본형 통합: Stage A 모듈 1 + 버스 감쇠기 + 조정형 Stage C & E23, E1 & 통과 3/3 \\
TB17 & R2-W460 통합 강등: TB13 + Stage A·조정형 Stage C 회로 & E23, E17 & 통과 2/2(강등 중 48 V ±0.03 \%) \\
TB18 & Stage A 경부하 위상 셰딩: 모듈 1 + R2 버스 + 충전기 전류원 & E24 & 통과 5/5 \\
TB19 & TB17 + 위상 셰딩 게이트 & E24, E17 & 통과 2/2 \\
TB20 & 충전 중 정지: Stage A 4모듈 + 충전기 + V\textsubscript{A} 클램프 + 버스 + 조정형 Stage C & E26 & 3/4(충전기 램프 사례 턴오프 1.01 A > 1.0 A) \\
\end{longtable}
}

\section{모델 간 차이와 해석상 유의점}

\begin{itemize}
\item \textbf{게이트 에지 양자화.} 게이트 에지가 적분 간격에 맞춰진다. 이 양자화가 정상상태 리밋 사이클 리플의 원인이었고(TB04·TB04B), 0.625 ns 이하에서 기준을 만족한다. TB06–TB20은 2.5 ns(TB06은 5 ns)로 과도와 한계를 판정했다.
\item \textbf{Stage A 등가.} TB08–TB13의 V\textsubscript{A}는 이상 전원(V\textsubscript{bat}/2 + 13.2 mΩ + 1.5 µH)이다. TB16–TB20은 이를 Stage A 회로로 바꿔 등가가 맞음을 확인했다. 정적 등가의 유효 범위는 V\textsubscript{A} 노드 공진(주기 약 30 µs)보다 느린 외란이다.
\item \textbf{Stage A 모듈 수.} TB16–TB19는 모듈 하나(2위상)로 계산했다. 4모듈 8위상은 TB07·TB20이다. 모듈 셰딩과 위상 셰딩의 조합은 회로로 확인하지 않았다.
\item \textbf{바이패스와 소프트 투입.} 바이패스·제동 저항 루프의 기생 인덕턴스와 바이패스 셀 출력용량은 넣지 않았다(투입 펄스는 상한 쪽 값). 소프트 투입은 가변 저항으로 모델링해 게이트 전압·채널 저항 관계와 소자 온도는 모델 밖이다.
\item \textbf{시간 축 단축.} 400 V 부하 분리·감발은 시간 지연 뒤 정전력 부하가 줄어드는 것으로 모델링했고, TB09–TB19는 그 시간을 3 ms로 줄였다(평균모델 기본 20 ms). TB20은 충전기 정지 시간을 3 ms로 줄였다(규격 30 ms, 흡수 에너지는 E26).
\item \textbf{충전기·부하.} 충전기는 이상 전류원, 400 V 부하는 정전력 등가다. 충전기의 출력 커패시터·전압 제한 루프·케이블 인덕턴스와 압축기·인버터의 스위칭 맥동은 넣지 않았다.
\item \textbf{해석기.} 후진 오일러는 공진망에 수치 감쇠를 만든다(TB05). 조정형 Stage C의 수 MHz 레그 공진도 감쇠되지만, 가변 간격 해석과의 대조(TB14 사례 A)에서 10 µs 평균의 차는 없었다. DAB 손실은 회로 저항분만 남는다.
\end{itemize}
